% miniVivisection manual -- the look: fonts, colors, headings, boxes, keys, code.
% Compiled with LuaLaTeX (fontspec, system fonts for symbols): tools/make_manual_latex.py.

% ---------------------------------------------------------------- fonts
\usepackage{fontspec}
\setmainfont{SourceSerif4}[
  UprightFont = *-Regular, ItalicFont = *-RegularIt, BoldFont = *-Semibold, BoldItalicFont = *-SemiboldIt,
  Extension = .otf, Numbers = Lining]
\setsansfont{SourceSans3}[
  UprightFont = *-Regular, ItalicFont = *-RegularIt, BoldFont = *-Semibold, BoldItalicFont = *-SemiboldIt,
  Extension = .otf]
\setmonofont{SourceCodePro}[
  UprightFont = *-Regular, ItalicFont = *-RegularIt, BoldFont = *-Semibold,
  Extension = .otf, Scale = 0.86]
\newfontfamily\mvhead{SourceSans3}[
  UprightFont = *-Semibold, BoldFont = *-Bold, ItalicFont = *-SemiboldIt, Extension = .otf, Ligatures = TeX]
\newfontfamily\mvlight{SourceSans3}[UprightFont = *-Light, Extension = .otf, Ligatures = TeX]
\newfontfamily\symbolfont{Segoe UI Symbol}
% The V° mark, in the font of the V° button on Grasshopper's toolbar.
\newfontfamily\vlogo{Segoe UI Bold}
% Outlined text: stroked, not filled (PDF text render mode 1), the stroke #1 bp wide.
\providecommand{\outlined}[2]{\pdfextension literal{1 Tr #1 w}#2\pdfextension literal{0 Tr}}
\usepackage{newunicodechar}
\newunicodechar{◆}{{\symbolfont ◆}}
\newunicodechar{☐}{{\symbolfont ☐}}
\newunicodechar{☑}{{\symbolfont ☑}}
\newunicodechar{∣}{{\symbolfont ∣}}
\newunicodechar{✓}{{\symbolfont ✓}}
\newunicodechar{▸}{{\symbolfont ▸}}
\newunicodechar{⟶}{{\symbolfont ⟶}}
\usepackage{microtype}
\usepackage[strings]{underscore} % _ in text: code names such as gh_item

% ---------------------------------------------------------------- colors
\usepackage[table]{xcolor}
\definecolor{mvblue}{HTML}{1A4D8F}
\definecolor{mvbluelight}{HTML}{EAF0F8}
\definecolor{mvbluemid}{HTML}{C9D6EA}
\definecolor{mvorange}{HTML}{D9661F}
\definecolor{mvorangelight}{HTML}{FCEFE6}
\definecolor{mvgreen}{HTML}{2E7D4F}
\definecolor{mvgreenlight}{HTML}{E8F4EC}
\definecolor{mvred}{HTML}{B3261E}
\definecolor{mvredlight}{HTML}{FBEAE9}
\definecolor{mvgray}{HTML}{5F6B7A}
\definecolor{mvgraylight}{HTML}{F2F4F7}
\definecolor{mvink}{HTML}{1D1D1F}
\definecolor{codebg}{HTML}{F6F8FA}
\definecolor{codekw}{HTML}{0000FF}
\definecolor{codestr}{HTML}{A31515}
\definecolor{codecom}{HTML}{008000}
\definecolor{codenum}{HTML}{098658}
\definecolor{codefn}{HTML}{795E26}
\color{mvink}

% ---------------------------------------------------------------- page
\usepackage[a4paper, top=24mm, bottom=26mm, inner=24mm, outer=24mm, headsep=9mm, footskip=12mm]{geometry}
\setlength{\parindent}{0pt}
\setlength{\parskip}{0.55em plus 0.2em minus 0.1em}
\linespread{1.07}
\usepackage{scrlayer-scrpage}
\pagestyle{scrheadings}
\clearpairofpagestyles
\automark[section]{chapter}
\ihead{\sffamily\small\color{mvgray}miniVivisection \,·\, Manual}
\ohead{\sffamily\small\color{mvgray}\leftmark}
\ofoot*{\sffamily\small\color{mvblue}\thepage}
\setkomafont{pageheadfoot}{\normalfont}
\renewcommand*{\chaptermarkformat}{\thechapter\enskip}

% ---------------------------------------------------------------- headings
\setkomafont{disposition}{\mvhead\color{mvblue}}
\setkomafont{chapter}{\mvhead\Huge\color{mvblue}}
\setkomafont{section}{\mvhead\Large\color{mvblue}}
\setkomafont{subsection}{\mvhead\large\color{mvink}}
\setkomafont{subsubsection}{\mvhead\normalsize\color{mvink}}
\setkomafont{paragraph}{\mvhead\normalsize\color{mvblue}}
\setkomafont{part}{\mvhead\Huge\color{mvblue}}
\setkomafont{partnumber}{\mvlight\Large\color{mvgray}}
\RedeclareSectionCommand[beforeskip=-1.6em plus -0.4em, afterskip=0.6em]{section}
\RedeclareSectionCommand[beforeskip=-1.2em plus -0.3em, afterskip=0.35em]{subsection}
\RedeclareSectionCommand[beforeskip=-0.9em plus -0.2em, afterskip=0.25em]{subsubsection}
\RedeclareSectionCommand[beforeskip=0.7em plus 0.2em, afterskip=-0.6em]{paragraph}
\RedeclareSectionCommand[beforeskip=0pt, afterskip=2.2em]{chapter}
\renewcommand*{\chapterformat}{%
  \begin{tikzpicture}[baseline=(n.base)]
    \node[fill=mvblue, text=white, inner xsep=7pt, inner ysep=5pt, rounded corners=3pt, font=\mvhead\huge] (n) {\thechapter};
  \end{tikzpicture}\hspace{0.6em}}
\renewcommand*{\chapterlinesformat}[3]{%
  \parbox{\linewidth}{\raggedright#2#3\par\vspace{0.35em}{\color{mvbluemid}\rule{\linewidth}{1.2pt}}}}
\renewcommand*{\sectionformat}{{\color{mvbluemid}\thesection}\hspace{0.7em}}

% ---------------------------------------------------------------- lists, tables, figures
\usepackage{enumitem}
\setlist{itemsep=0.2em, topsep=0.3em, parsep=0.15em, leftmargin=1.4em}
\setlist[itemize,1]{label={\color{mvblue}\small\textbullet}}
\setlist[itemize,2]{label={\color{mvgray}\footnotesize–}}
\setlist[description]{font=\sffamily\bfseries\color{mvink}, itemsep=0.25em, leftmargin=1.4em}
\usepackage{booktabs, tabularx, longtable, xltabular, array, multirow}
\renewcommand{\arraystretch}{1.22}
\newcolumntype{L}{>{\raggedright\arraybackslash}X}
\newcolumntype{P}[1]{>{\raggedright\arraybackslash}p{#1}}
\setlength{\aboverulesep}{0.25ex}
\setlength{\belowrulesep}{0.45ex}
\arrayrulecolor{mvbluemid}
\usepackage{graphicx}
% The screenshots as high-quality JPEGs, made from docs/screenshots/*.png by tools/make_manual_latex.py.
\graphicspath{{generated/jpg/}}
\usepackage[font={small,sf}, labelfont={bf,color=mvblue}, justification=raggedright, singlelinecheck=false, skip=6pt]{caption}
\usepackage{float}

% ---------------------------------------------------------------- tikz, boxes
\usepackage{tikz}
\usetikzlibrary{arrows.meta, positioning, shapes.geometric, calc, fit, backgrounds}
% TikZ's ragged alignment sets \spaceskip for Computer Modern -- far too wide in Source Sans: the font's own spaces.
\tikzset{every node/.append style={execute at begin node={\spaceskip=0pt\relax\xspaceskip=0pt\relax}}}
\usepackage[most]{tcolorbox}
\tcbset{
  mvbox/.style={enhanced, breakable, boxrule=0pt, arc=2pt, left=10pt, right=10pt, top=6pt, bottom=6pt,
    before skip=0.8em, after skip=0.8em, fonttitle=\sffamily\bfseries\small, coltitle=#1, attach title to upper={\enskip}},
}
\newtcolorbox{note}[1][Note]{mvbox=mvblue, colback=mvbluelight, borderline west={3pt}{0pt}{mvblue}, title={#1}}
\newtcolorbox{tip}[1][Tip]{mvbox=mvgreen, colback=mvgreenlight, borderline west={3pt}{0pt}{mvgreen}, title={#1}}
\newtcolorbox{caution}[1][Careful]{mvbox=mvorange, colback=mvorangelight, borderline west={3pt}{0pt}{mvorange}, title={#1}}
\newtcolorbox{lesson}[1]{mvbox=mvgray, colback=mvgraylight, borderline west={3pt}{0pt}{mvgray}, title={Try it: #1}}
% A concept: a titled box of its own.
\newtcolorbox{concept}[1]{enhanced, breakable, colback=white, colframe=mvbluemid, boxrule=0.8pt, arc=3pt,
  left=10pt, right=10pt, top=8pt, bottom=8pt, before skip=1em, after skip=1em,
  title={#1}, fonttitle=\mvhead\normalsize, coltitle=mvblue, colbacktitle=mvbluelight,
  attach boxed title to top left={xshift=8pt, yshift=-8pt}, boxed title style={boxrule=0.8pt, colframe=mvbluemid, arc=2pt}}

% ---------------------------------------------------------------- inline markup
% \vn{Watch} -> V°Watch, the plug-in's V names.
\newcommand{\vn}[1]{\mbox{\sffamily\bfseries V°#1}}
% \ui{Edit in Rhino}: a label, button or menu item as the user sees it.
\newcommand{\ui}[1]{{\sffamily\bfseries #1}}
% \menu{A}{B}: a path through menus or settings groups.
\newcommand{\menu}[1]{{\sffamily\bfseries\def\>{\,{\color{mvgray}\small▸}\,}#1}}
% \setting{V°Settings > Debugger > Stop once per line}
\newcommand{\setting}[1]{{\sffamily\itshape\color{mvblue}#1}}
% \code{...}: code in the text.
\newcommand{\code}[1]{\texttt{\color{mvink}#1}}
\setlength{\emergencystretch}{2.5em}
\newcommand{\codebreak}[1]{\texttt{\color{mvink}#1}}
% \bs: a backslash in a path, where the line may break.
\newcommand{\bs}{\textbackslash\allowbreak}
% \file{...}: file and folder names.
\newcommand{\file}[1]{\texttt{\color{mvgray}#1}}
% \key{Ctrl}: a key cap; \keys{Ctrl+Alt+B} for a chord.
\newcommand{\key}[1]{\tikz[baseline=(k.base)]\node[draw=mvgray!70, fill=mvgraylight, rounded corners=1.5pt, inner xsep=2.6pt, inner ysep=1.4pt, line width=0.45pt, font=\sffamily\scriptsize] (k) {\strut #1};}
\ExplSyntaxOn
\NewDocumentCommand{\keys}{m}{
  \seq_set_split:Nnn \l_tmpa_seq {+} {#1}
  \bool_set_true:N \l_tmpa_bool
  \seq_map_inline:Nn \l_tmpa_seq {
    \bool_if:NF \l_tmpa_bool { \mbox{\scriptsize\,+\,} }
    \bool_set_false:N \l_tmpa_bool
    \key{##1}
  }
}
\ExplSyntaxOff
% A reference to a section, shown as "§ 5.2" with a link.
\newcommand{\see}[1]{\mbox{\hyperref[#1]{§\,\ref*{#1}}}}
% A lesson's file.
\newcommand{\lessonfile}[1]{\texttt{\color{mvgreen}#1}}

% ---------------------------------------------------------------- code listings
\usepackage{listings}
\lstdefinestyle{mv}{
  basicstyle=\ttfamily\small, columns=fullflexible, keepspaces=true,
  backgroundcolor=\color{codebg}, frame=single, framerule=0pt, rulecolor=\color{codebg},
  framesep=6pt, xleftmargin=6pt, xrightmargin=6pt, aboveskip=0.8em, belowskip=0.8em,
  keywordstyle=\color{codekw}, stringstyle=\color{codestr}, commentstyle=\color{codecom}\itshape,
  showstringspaces=false, upquote=true, breaklines=true, breakatwhitespace=true,
  literate={°}{{°}}1 {→}{{$\rightarrow$}}1 {…}{{\dots}}1 {◆}{{\symbolfont ◆}}1,
}
\lstdefinelanguage{PythonMV}[]{Python}{morekeywords={True,False,None,as,with,async,await,nonlocal}}
\lstset{style=mv, language=PythonMV}
\lstdefinestyle{plain}{style=mv, language={}}

% ---------------------------------------------------------------- links
\usepackage[hidelinks, bookmarksnumbered, bookmarksopen, bookmarksopenlevel=1, pdfdisplaydoctitle]{hyperref}
\usepackage{bookmark}
\hypersetup{colorlinks=true, linkcolor=mvblue, urlcolor=mvblue, citecolor=mvblue}
\usepackage[capitalise, noabbrev]{cleveref}

% ---------------------------------------------------------------- table of contents
\setkomafont{chapterentry}{\mvhead\color{mvblue}}
\setkomafont{partentry}{\mvhead\large\color{mvink}}
\RedeclareSectionCommand[tocentryformat=\mvhead\color{mvblue}, tocpagenumberformat=\mvhead\color{mvblue}]{chapter}
\setcounter{tocdepth}{1}

% A figure with one of the screenshots (960 x 600): \shot{file}{caption}{label}
\newcommand{\shot}[3]{%
  \begin{figure}[htbp]
    \centering
    \tikz\node[inner sep=0pt, draw=mvbluemid, line width=0.6pt]{\includegraphics[width=0.92\linewidth]{#1}};
    \caption{#2}\label{#3}
  \end{figure}}

% A feature list's table: \begin{features}{Heading} \feat{name}{what}{section} ... \end{features}
\usepackage{needspace}
\newenvironment{features}[1]{%
  \par\needspace{7\baselineskip}\smallskip{\mvhead\color{mvblue}#1}\par\nopagebreak\vspace{-0.4em}
  \xltabular{\linewidth}{@{}P{0.27\linewidth}L>{\raggedleft\arraybackslash}p{0.07\linewidth}@{}}
  \toprule}{\bottomrule\endxltabular}
\newcommand{\feat}[3]{{\sffamily\bfseries #1} & #2 & \see{#3}\\}
